%% kapa: exam=1 round=35 year=2024 period=1 form=A subject=civil_law subject_ko=민법 seq=37 number=37
%% source: Q-Net 감정평가사 기출문제 (cst00309) artlSeq=5213229 file=raw/1008/5213229/1교시 문제지 원본.pdf
%% conversion: kordoc 4.24.51 -> tools/gen1.py (auto); review=proofread-sonnet; flags=none
\begin{question}{37}
질권에 관한 설명으로 옳지 않은 것은?
\begin{choices}
  \item 질물보다 다른 재산이 먼저 경매된 경우, 질권자는 그 매각대금으로부터 배당을 받을 수 없다.
  \item 질권자가 채권 일부를 변제받았더라도 질물 전부에 대하여 그 권리를 행사할 수 있다.
  \item 질물이 멸실된 경우에도 그로 인하여 질권설정자가 받을 금전을 압류하면 질권의 효력이 그 금전에 미친다.
  \item 정당한 이유있는 때에는 질권자는 채무자 및 질권설정자에게 통지하고 감정자의 평가에 의하여 질물로 직접 변제에 충당할 것을 법원에 청구할 수 있다.
  \item 질권자는 그 권리의 범위 내에서 자기의 책임으로 질물을 전질할 수 있다.
\end{choices}
\end{question}
\answer{1}
